\documentclass{article}
\usepackage[T1]{fontenc}
\usepackage[utf8]{inputenc}
\usepackage{cancel}
\usepackage{amsmath}
\usepackage{amssymb}

\title{High School}
\author{Wiktoria Borek}
\date{}

\begin{document}

\maketitle

\section{Reasoning with Real Numbers}
\subsection{Exercise}
Prove that $2+2=5$.
\subsubsection*{Solution}
$
  0=0 \rightarrow 20 -16 -4 = 25 -20 -5 \\
  4 (5-4-1) = 5 (5-4-1) \\
  4 = 5 \\
  2 + 2 = 5 \quad \blacksquare \\
  \\
  \uparrow \text{This is a joke, not a real proof.} \\
  \text{And that's because the number $5-4-1$ is equal to $0$, and we cannot divide by $0$.}
$

\section{Logarithms}
\subsection{Exercise}
Prove that if $x>1$, then: \\
$2\log_2{(x^2-1)} - 3\log_2{(x-1)} = \log_2{\frac{x^2+2x+1}{x-1}}$
\subsubsection*{Solution}
$
  D: x>1 \\
    \\
    L = 2\log_2{(x^2-1)} - 3\log_2{(x-1)} \\
    R = \log_2{\frac{x^2+2x+1}{x-1}} \\
    \\
    L = \log_2{(x^2-1)^2} - \log_2{(x-1)^3} = \log_2{\frac{(x^2-1)^2}{(x-1)^3}} = \log_2{\frac{\cancel{(x-1)^2}(x+1)^2}{\cancel{(x-1)^2}(x-1)}} = \\
    = \log_2{\frac{x^2+2x+1}{x-1}} = R \quad \blacksquare
$
\subsection{Exercise}
Prove that $\log_{17}{19} : \log_{18}{19} = \log_{16}{18} \cdot \log_{17}{16}$.
\subsubsection*{Solution}
$
  \log_{17}{19} \cdot \log_{19}{18} = \log_{17}{16} \cdot \log_{16}{18} \\
  \log_{17}{\cancel{19}} \cdot \log_{\cancel{19}}{18} = \log_{17}{\cancel{16}} \cdot \log_{\cancel{16}}{18} \\
  \log_{17}{18} = \log_{17}{18} \\
  L = R \quad \blacksquare
$
\subsection{Exercise}
Given the numbers $a=\log_3{2}$ and $b=\log_{2^{2024}}{3^{1012}} + \log_{\frac{1}{3}}{54}$. Express $b$ in terms of $a$.
\subsubsection*{Solution}
$
  b = \frac{1}{2} \cdot \log_2{3} - \log_3{(27 \cdot 2)} = \frac{1}{2} \cdot \frac{1}{a} - (\log_3{3^3} + a) = \frac{1}{2a} - 3 - a = \frac{1}{2a} - \frac{6a}{2a} - \frac{2a^2}{2a} = \\
    = - \frac{2a^2 +6a - 1}{2a}
$
\subsection{Exercise}
Given that $\frac{1}{\log_{p-q}{2}} = 3$ and $(\log_{pq}{3})^{-1}=2$. Prove that the number \\
$\sqrt{p^2 - 1 + q^2}$ is rational.
\subsubsection*{Solution}
$
  \log_{2}{p-q} = 3 \quad \land \quad \log_{3}{pq} = 2 \rightarrow pq = 3^2 = 9 \\
  p-q = 2^3 \; / \; ^2 \\
  (p-q){^2} = 64 \\
  p^2 + q^2 - 2pq = 64 \\
  p^2 + q^2 = 82 \\
  \\
  \sqrt{p^2 - 1 + q^2} = \sqrt{p^2 + q^2 - 1} = \sqrt{82 - 1} = \sqrt{81} = 9 \in \mathbb{Q} \quad \blacksquare
$
\subsection{Exercise}
Given that $\log_a{m} = \sqrt{2}$. Evaluate $3\log_{\sqrt{m}}{am^2} - \log_{am}{\frac{a^2}{m}}$.
\subsubsection*{Solution}
$
  3 \cdot 2 \log_{m}{am^2} - \frac{\log_a{\frac{a^2}{m}}}{\log_a{am}} = 6 \cdot \log_m{a} + 12 - \frac{\log_a{a^2} - \log_a{m}}{\log_a{a} + \log_a{m}} = \\
   = \frac{6}{\sqrt{2}} + 12 - \frac{2 - \sqrt{2}}{1 + \sqrt{2}} = \frac{6\sqrt{2}}{2} + 12 - \frac{(2 - \sqrt{2})(1 - \sqrt{2})}{(1 + \sqrt{2})(1 - \sqrt{2})} = \\
   = \frac{24 + 6\sqrt{2}}{2} - \frac{2 - 2\sqrt{2} - \sqrt{2} + 2}{1 - 2} = \frac{24 + 6\sqrt{2}}{2} + 4 - 3\sqrt{2} = 16
$
\subsection{Exercise}
Prove that for $a, b, c > 1$ and $n \in \mathbb{N}$ the number $\frac{\log_a{b^3} \cdot \log_{b^2}{c}}{\log_{\sqrt[3]{a^n}}{\sqrt[4]{c}}}$ is even.
\subsubsection*{Solution}
$
 \frac{\frac{3}{2} \log_a{b} \cdot \log_b{c}}{\frac{3}{4}n \log_a{c}} = \frac{\frac{3}{2} \cancel{\log_a{c}}}{\frac{3}{4}n \cancel{\log_a{c}}} = \frac{3}{2} \cdot \frac{4}{3} \cdot n = 2 \cdot n \quad \blacksquare
$
\end{document}
